\documentclass[11pt]{article}
\usepackage[a4paper,margin=21mm]{geometry}
\usepackage[no-math]{fontspec}
\setmainfont{lmroman10-regular.otf}[BoldFont=lmroman10-bold.otf,ItalicFont=lmroman10-italic.otf,BoldItalicFont=lmroman10-bolditalic.otf]
\usepackage{amsmath,amssymb,amsthm,mathtools,stmaryrd,mathrsfs}
\usepackage{csquotes,longtable,tabu}
\usepackage[shortlabels]{enumitem}
\usepackage{xurl}
\usepackage[unicode,hidelinks]{hyperref}
\usepackage[nameinlink]{cleveref}
\crefdefaultlabelformat{#2\textup{#1}#3}
\allowdisplaybreaks
\setlength\emergencystretch{3em}
\newlength{\normalparindent}
\setlength{\normalparindent}{\parindent}
\providecommand{\ie}{i.e.\ }
\providecommand{\oldbigcup}{\bigcup}
\theoremstyle{plain}
\newtheorem{theorem}{Theorem}[section]
\usepackage{aliascnt}
\newaliascnt{lemma}{theorem}
\newtheorem{lemma}[lemma]{Lemma}
\aliascntresetthe{lemma}
\crefname{lemma}{Lemma}{Lemmas}
\newtheorem{proposition}[theorem]{Proposition}
\newtheorem{corollary}[theorem]{Corollary}

\theoremstyle{definition}
\newtheorem{definition}[theorem]{Definition}
\newtheorem{remark}[theorem]{Remark}
\newtheorem{example}[theorem]{Example}

\newcommand{\dd}{\mathop{}\mathopen{}\mathrm{d}}
\newcommand{\C}{\mathbb{C}}
\newcommand{\N}{\mathbb{N}}
\newcommand{\Q}{\mathbb{Q}}
\newcommand{\R}{\mathbb{R}}
\newcommand{\Z}{\mathbb{Z}}
\newcommand{\Sph}{\mathbb{S}}
\newcommand{\Cspace}{\mathscr{C}}
\newcommand{\Mspace}{\mathscr{M}}
\newcommand{\ener}{\mathcal{E}}
\newcommand{\mass}{\boldsymbol{M}}
\newcommand{\mres}{\mathop{\hbox{\vrule height 6pt width .5pt depth 0pt \vrule height .5pt width 4pt depth 0pt}}\nolimits}
\newcommand{\id}{\mathrm{Id}}
\newcommand{\loc}{\mathrm{loc}}
\newcommand{\lsc}{\mathrm{lsc}}
\newcommand{\eps}{\varepsilon}
\newcommand{\uloc}{L^1_\mathrm{uloc}}
\newcommand{\st}{\;:\;}

\DeclarePairedDelimiter{\abs}{\lvert}{\rvert}
\DeclarePairedDelimiter{\norm}{\lVert}{\rVert}
\DeclarePairedDelimiterX{\intint}[2]{\llbracket}{\rrbracket}{#1,#2}
\DeclareMathOperator{\supp}{supp}
\DeclareMathOperator{\diam}{diam}
\DeclareMathOperator{\tang}{Tan}
\DeclareMathOperator*{\esssup}{ess\,sup}
\DeclareMathOperator{\hdm}{\mathcal{H}}
\DeclareMathOperator{\lbm}{\mathcal{L}}
\DeclareMathOperator{\dist}{dist}
\DeclareMathOperator{\reach}{reach}

\newcommand{\weakto}{\rightharpoonup}
\newcommand{\mweakto}{\xrightharpoonup{\Cspace_0'}}
\newcommand{\narrowto}{\xrightharpoonup{\Cspace_b'}}
\newcommand{\xto}[2][]{\xrightarrow[#2]{#1}}
\newcommand{\xmweakto}[1]{\xrightharpoonup[#1]{\Cspace_0'}}
\newcommand{\xweakto}[1]{\xrightharpoonup[#1]{}}
\newcommand{\xnarrowto}[1]{\xrightharpoonup[#1]{\Cspace_b'}}

\numberwithin{equation}{section}
\begin{document}
\begin{center}\Large Countable bubble decomposition with conservation of bubbling mass for bounded positive measures\end{center}
\noindent\textbf{Stable ID:} 2243\par
\noindent\textbf{Author:} Antonin Monteil; Paul Pegon\par
\noindent\textbf{Work:} Mass concentration in rescaled first order integral functionals\par
\noindent\textbf{Source:} \url{https://jep.centre-mersenne.org/articles/10.5802/jep.257/}\par
\noindent\textbf{Version:} Journal de l'Ecole polytechnique - Mathematiques 11 (2024),431--472; published2024-02-19; original Theorem 3.2, printed 448--452 / physical PDF 19--23; official native/PDF hashes pinned\par
\noindent\textbf{Rights:} CC BY4.0. \url{https://creativecommons.org/licenses/by/4.0/}. Independent typesetting; author mathematics preserved.\par
\noindent\textbf{Difficulty:} H2. Banach-Alaoglu compactness and Portmanteau inequalities only yield local vague limits and allow mass to escape; they do not provide one subsequence with countably many simultaneously separated profiles. The authors construct nested extractions, a diagonal family of expanding balls and quantitative mass errors, then control every translated limit of the remainder by a decreasing bubbling quantity. They finally recenter each bubble in the support of the original measure while preserving narrow convergence and the whole bubbling-mass sum, which is nonroutine beyond qualifying-examination measure compactness.\par
\noindent\textbf{Editorial scope:} One original published Theorem 3.2: all five conclusions (A)--(E), finite/countably infinite and vanishing cases, support-centered balls, narrow/vague distinctions and signed residual construction are retained. Bubbling mass conservation does not claim that the vanishing part has zero total mass. The full author six-step argument and the complete local vanishing characterization are included.\par
\medskip\hrule\medskip
% Exact originalnative source lines 118--132
\section*{Notation}
\begin{tabu} {X[2,c,m] X[12,l,p]}
\(B_r(x)\) & open ball of radius \(r\) centered at \(x\);\\
\(B_r\) & open ball $B_r(0)$;\\
\(\Mspace(\R^N)\) & set of finite signed Borel measures on \(\R^N\);\\
\(\Mspace_+(\R^N)\) & set of finite positive Borel measures on \(\R^N\);\\
$\Phi_\sharp\mu$ & pushforward of a measure $\mu\in\Mspace(\R^N)$ by a map $\Phi:\R^N\to\R^k$, defined as $A\mapsto \mu(\Phi^{-1}(A))$;\\
$\tau_x\mu$ & Borel measure $A \mapsto \mu(A-x)$ if \(\mu\in\Mspace(\R^N)\) and $x\in \R^N$;\\
$c_B \mu$ & Borel measure $\tau_{-x} (\mu \mres B)$ if $B$ is the ball $B_r(x)$;\\
$\mu_\ell \mweakto \mu$ & weak convergence of measures, \ie weak-$\star$ convergence in duality with the space $\Cspace_0(\R^N)$ of continuous functions vanishing at infinity;\\
$\mu_\ell \narrowto \mu$ & narrow convergence of measures, \ie weak-$\star$ convergence in duality with the space $\Cspace_b(\R^N)$ of continuous and bounded function;\\
$\Sigma$ & set of increasing maps $\sigma : \N \to \N$;\\
$\sigma_1 \preceq \sigma_2$ & $\sigma_1, \sigma_2 \in \Sigma$ are such that $\sigma_1(\intint{n}{+\infty}) \subseteq \sigma_2(\N)$ for some $n\in \N$ ;\\
$\pm$ & fixed to either $+$ or $-$ in the whole statement or proof, and $\mp = -(\pm)$.
\end{tabu}
\setcounter{section}{2}
% Exact originalnative source lines 690--690
\section{Lower bound for the energy and existence of optimal profiles}\label{sectionLowerBound}
% Exact originalnative source lines 694--704
\subsection{Profile decomposition by concentration-compactness}

We prove a profile decomposition theorem for bounded sequences of positive measures over $\R^N$, which is essentially equivalent to \cite[Th.\,1.5]{marisProfileDecompositionSequences2014} in the Euclidean case. We have added an extra information on mass conservation that will be useful, and provide a self-contained simple proof. We start with a definition.

\begin{definition}\label{DefinitionVanishing}
A sequence of positive measures $(\mu_n)_{n\in\N} \!\in\! \Mspace_+(\R^N)$ is \emph{vanishing}~if
\[\sup_{x\in \R^N} \mu_n(B_1(x)) \xto{n\to \infty} 0.\]
\end{definition}

Any bounded sequence of positive measures over $\R^N$ may be decomposed (up~to subsequence) into a countable collection of narrowly converging \enquote{bubbles} and a vanishing part, accounting for the total mass of the sequence, as stated in the following theorem.

% Exact originalnative source lines 705--718
\begin{theorem}\label{profileDecomposition}
For every bounded sequence \((\mu_n)_{n\in\N}\) of positive Borel measures on~\(\R^N\), there exists a subsequence \((\mu_n)_{n \in \sigma(\N)}\), $\sigma \in \Sigma$, a non-decreasing sequence of integers $(k_n)_{n \in \sigma(\N)}$ converging to some $k \in \N \cup \{+\infty\}$, a sequence of non-trivial positive Borel measures $(\mu^i)_{0\le i<k}$, and for every $n \in \sigma(\N)$, a collection of balls $(B^i_n)_{0 \leq i < k_n}$ centered at points of $\supp \mu_n$ such that, writing for all $n\in \sigma(\N)$,
\begin{equation}
\mu_n = \mu^b_n + \mu^v_n,\quad \text{where } \mu^b_n = \sum_{0\le i < k_n} \mu_n \mres B_n^i,
\end{equation}
\vspace*{-5pt}
\begin{enumerate}[\textnormal(\textup{\Alph*}\textnormal),leftmargin=1.3\normalparindent]
\item\label{bubbleEmergence} bubbles emerge: $(c_{B_{n}^i} \mu_{n})_{n\in\sigma(\N)}\xnarrowto{n\to\infty}\mu^i$ for every $i<k$,\footnote{Recall that $c_B \mu = (x\mapsto x-y)_\sharp(\mu \mres B)$ if $B = B_r(y)$ and $\mu \in \Mspace(\R^N)$.}
\item\label{bubbleSplitting} bubbles split: $\min_{0\le i<j< k_{n}}\dist(B_{n}^i,B_{n}^j) \xto{n\to\infty} +\infty$,
\item\label{bubbleDivergence} bubbles diverge: $ \min_{0\le i< k_n}\diam(B_{n}^i) \xto{n\to\infty} +\infty$,
\item\label{bubblingMassConservation} the bubbling mass is conserved: $\norm{\mu^b_{n}} \xto{\ell \to \infty} \sum_{0\le i<k} \norm{\mu^i}$,
\item\label{remainingVanishing} the remaining part is vanishing: $\sup_{x\in \R^N} \mu^v_{n}(B_1(x)) \xto{n\to\infty} 0$.
\end{enumerate}
\end{theorem}
% Exact originalnative source lines 720--740
Before proving \Cref{profileDecomposition}, we introduce the \enquote{bubbling} function of a sequence of finite \emph{signed} measures \((\mu_n)_{n\in\N}\):\vspace*{-3pt}
\begin{equation}
m((\mu_n)_{n\in\N})\coloneqq\sup\Big\{ \norm{\mu} \st (\tau_{-x_{\sigma(\ell)}}\mu_{\sigma(\ell)})_{\ell\in\N} \mweakto \mu,\, \sigma \in \Sigma,\, x_{\sigma(\ell)} \in \R^N\, (\forall \ell) \Big\}.
\end{equation}
Although we will use this function on signed measures, we will start from a sequence of positive measures and use the following characterization of vanishing sequences, which holds only in the case of positive measures:
\begin{lemma}\label{concentrationProperties}
A sequence \((\mu_n)_{n\in\N}\) of finite positive measures over $\R^N$ is vanishing if and only if \(m((\mu_n)_{n\in\N})=0\).
\end{lemma}
\begin{proof}
Assume that \((\mu_n)_{n\in\N}\) is vanishing and that \((\tau_{-x_{\sigma(\ell)}}\mu_{\sigma(\ell)})_{\ell\in\N} \mweakto \mu\) for some \(\sigma \in \Sigma\) and some sequence of points \((x_{\sigma(\ell)})_{\ell\in\N}\). Then, for every \(x\in\R^N\),\vspace*{-3pt}
\[
\mu(B_1(x))\le\liminf_{\ell\to\infty}\tau_{-x_{\sigma(\ell)}}\mu_{\sigma(\ell)}(B_1(x))=\liminf_{\ell\to\infty}\mu_{\sigma(\ell)}(B_1(x+x_{\sigma(\ell)}))=0,
\]
\ie \(\mu=0\) and thus \(m((\mu_\ell)_{\ell\in\N})=0\).

Conversely, if \((\mu_n)_{n\in\N}\) is not vanishing, then there exists \(\eps>0\), \(\sigma\in\Sigma\) and a sequence of points \((x_n)_{n\in\sigma(\N)}\) in \(\R^N\) such that \(\mu_n(B_1(x_n))\ge\eps\) for every \(n\in\sigma(\N)\). Up to further extraction, one can assume that \((\tau_{-x_{\sigma(\ell)}}\mu_{\sigma(\ell)})_{\ell\in\N} \mweakto\mu\in\Mspace(\R^N)\). We~have\vspace*{-3pt}
\[
\mu(\bar B_1(0))\ge\limsup_{\ell\to\infty}\tau_{-x_{\sigma(\ell)}}\mu_{\sigma(\ell)}(\bar B_1(0))=\limsup_{\ell\to\infty}\mu_{\sigma(\ell)}(\bar B_1(x_{\sigma(\ell)}))\ge\eps>0,
\]
which entails \(m((\mu_\ell)_{\ell\in\N})\ge\eps>0\).
\end{proof}
% Exact originalnative source lines 742--857
\begin{proof}[Proof of \Cref{profileDecomposition}]
If \((\mu_n)_{n\in\N}\) is vanishing, then we take $\sigma = \id$ and $k=0$, so~that $\mu_{\sigma(\ell)} = \mu_\ell = \mu^v_\ell$, \labelcref{bubbleEmergence,bubbleSplitting,bubbleDivergence,bubblingMassConservation} are empty statements and \ref{remainingVanishing} is satisfied since $(\mu_n)_{n\in\N}$ is vanishing. Assume on the contrary that \((\mu_n)_{n\in\N}\) is not vanishing. We shall construct the bubbles by induction and prove their properties in several steps.

\subsubsection*{Step 1: construction of bubbles centers} At first step (step $0$), since \(m((\mu_n)_{n\in\N})>0\), there exists $\sigma_0 \in \Sigma$ and a sequence of points \((x^0_{n})_{n \in \sigma_0(\N)}\), such that
\begin{equation}
(\tau_{-x^0_{n}} \mu_n)_{n \in \sigma_0(\N)}\mweakto\mu^0\in\Mspace(\R^N)\quad\text{with}\quad\norm{\mu^0} \geq \frac 12 m((\mu_n)_{n\in\N}).
\end{equation}
We then set \(\mu^0_n\coloneqq\mu_n-\tau_{x^0_n}\mu^0\) and we continue by induction, starting from the sequence \((\mu^0_n)_{n\in\sigma_0(\N)}\). More precisely, assume that for a fixed step $k-1 \in \N$, for every \(i\in\N\) such that \(0\le i\le k-1\), we have built $\mu^i \in \Mspace(\R^N)$, $\sigma_i \in \Sigma$, points $(x^i_n)_{n\in \sigma_i(\N)}$ and sequences $(\mu^i_{n})_{n\in \sigma_i(\N)} \in \Mspace(\R^N)$ such that for every $i$,
\begin{gather}
\sigma_i \preceq \sigma_{i-1},\label{bubbleInduction1}\\
\mu^i_{n} = \mu_n -\sum_{0\leq j \leq i} \tau_{x^j_n}\mu^j, \quad (\forall n \in \sigma_i(\N)),\label{bubbleInduction2}\\
(\tau_{-x^i_n} \mu^{i-1}_n)_{n\in\sigma_i(\N)} \mweakto \mu^i,\label{bubbleInduction3}\\
\norm{\mu^i} \geq \frac 12 m((\mu^i_n)_{n \in\sigma_i(\N)})>0,\label{bubbleInduction4}
\end{gather}
where $\sigma_{-1} \coloneqq \id, (\mu_n^{-1}) \coloneqq (\mu_n)$. If $m((\mu^{k-1}_n)_{n\in \sigma_{k-1}(\N)}) = 0$, we stop; otherwise, we proceed to the next step $k$ to build $\sigma_k, \mu^k, (x_n^k)_{n\in\sigma_k(\N)}, (\mu_n^k)$ as we did at step $k=0$, starting with $(\mu^{k-1}_n)_{n\in\sigma_{k-1}(\N)}$. Either the induction stops at some step \(k-1\in\N\) for which \(m((\mu^{k-1}_{n})_{n\in \sigma_{k-1}(\N)})=0\) or the previous objects are defined for every \(i\in\N\), in~which case we let \(k \coloneqq +\infty\).

\subsubsection*{Step 2: splitting of bubbles centers} We prove that
\begin{equation}\label{splittingWeak}
\lim_{\sigma_i(\N) \ni n\to\infty} \dist(x^i_{n},x^j_{n}) =+\infty \quad\text{for every \(i,j\in\N\) with \(0\le j<i<k\)}.
\end{equation}
Indeed, assume by contradiction that there is a first index \(i<k\) such that for some \(j_0<i\), \((\dist(x^i_{n},x^{j_0}_{n}))_{n\in\sigma_i(\N)}\) is not divergent. In particular, there exists \(\sigma\preceq\sigma_i\) such that \((x^i_{n}-x^{j_0}_{n})_{n\in\sigma(\N)}\to x\in\R^N\). Moreover, \((\dist(x^i_n,x^j_n))_{n\in\sigma_i(\N)}\to\infty\), for every \(j<i\), \(j\neq j_0\) by minimality of $i$ and the triangle inequality $\dist(x_n^j,x_n^{j_0}) \leq \dist(x_n^j,x_n^i) + \dist(x_n^i,x_n^{j_0})$. Notice by \labelcref{bubbleInduction2} that for every $n\in \sigma(\N)$,
\begin{gather*}
\mu_n^{i-1} = \mu_n^{j_0-1} - \tau_{x_n^{j_0}} \mu^{j_0} - \sum_{j_0 < j < i} \tau_{x_n^j} \mu^j,\\
\shortintertext{hence taking the translation $\tau_{-x_n^i}$,}
\tau_{-x_n^i} \mu_n^{i-1} = \tau_{x_n^{j_0}-x_n^i} (\tau_{-x_n^{j_0}} \mu_n^{j_0-1}-\mu^{j_0}) - \sum_{j_0 < j < i} \tau_{x_n^j-x_n^i} \mu^j,\\
\shortintertext{and passing to the weak limit, knowing that $x_n^{j_0}-x_n^i \to -x$ and $\dist(x_n^j,x_n^i) \to +\infty$ for $j_0 < j <i$,}
\mu^i = \tau_{-x}(\mu^{j_0}-\mu^{j_0}) - \sum_{j_0 < j < i} 0 = 0.
\end{gather*}
This contradicts the fact that \((\tau_{-x^i_n}\mu^{i-1}_n)_{n\in\sigma(\N)}\mweakto\mu^i\neq 0\) and proves \labelcref{splittingWeak}.

\subsubsection*{Step 3: weak convergence of bubbles}
From \labelcref{bubbleInduction2} we get\vspace*{-3pt}
\begin{equation}\label{bubbleWeak}
\tau_{-x^i_n}\mu^{i-1}_n =\tau_{-x^i_n}\mu_n-\sum_{0\le j<i}\tau_{-x^i_n+x^j_n}\mu^j,
\end{equation}
and by \labelcref{splittingWeak}, the sum converges weakly to $0$, and so\vspace*{-3pt}\enlargethispage{\baselineskip}
\begin{equation}\label{bubbleWeakConvergence}
(\tau_{-x^i_n}\mu_n)_{n\in\sigma_i(\N)}\mweakto \mu^i\quad\text{for every \(i\in\N\) with \(i<k\)}.
\end{equation}

\subsubsection*{Step 4: construction of the bubbles with mass conservation}
We now construct the extraction \(\sigma\in\Sigma\) that we need by induction: we set \(\sigma(0)=0\) and, assuming that \(\sigma(0)<\dots<\sigma(\ell-1)\), with \(\ell\in\N^\ast\), have been constructed, we set \(\sigma(\ell)\coloneqq n\) with \(n\in\sigma_{\ell \wedge k -1}(\N)\) large enough so that $n > \sigma(\ell-1)$ and for every \(i< \ell \wedge k\),\vspace*{-3pt}
\begin{gather}
\mu_n(B_\ell(x^i_n))\le\norm{\mu^i}+2^{-\ell},\label{inductionBubbles1}\\[-3pt]
\shortintertext{and}
\min_{0\le j<i}\dist(x^i_n,x^j_n)\ge 4\ell.\label{inductionBubbles2}
\end{gather}
Such an \(n\) exists by \labelcref{splittingWeak} and \labelcref{bubbleWeakConvergence}, noticing that $\mu_n(B_\ell(x_n^i)) = (\tau_{-x_n^i} \mu_n)(B_\ell)$. Then for each \(n=\sigma(\ell)\), \(\ell\in\N\), we set \(k_n=\ell \wedge k\), and for each \(i\in\{0,\dots,k_n-1\}\),\vspace*{-3pt}
\[
B^i_n\coloneqq B_\ell(x^i_n).
\]
Finally, for every \(n\in\sigma(\N)\), we decompose $\mu_n$ as expected:\vspace*{-3pt}
\[
\mu_n = \mu^b_n + \mu^v_n,\quad \text{where } \mu^b_n = \sum_{0\le i < k_n} \mu_n \mres B_n^i.
\]

Let us check the four first items \labelcref{bubbleEmergence}--\labelcref{bubblingMassConservation}. Notice that \labelcref{bubbleDivergence} is fulfilled because $\diam(B_{\sigma(\ell)}^i) = \ell \to +\infty$ as $\ell\to\infty$, and \labelcref{bubbleSplitting} because of \labelcref{inductionBubbles2}. Since for every $i<k$, $\lim_{\sigma(\N)\ni n \to\infty}\diam(B_n^i) = +\infty$ and $c_{B_n^i} \mu_n = (\tau_{-x_n^i} (\mu_n\mres B_n^i))$ for every $n \in \sigma_i(\N)$, $(c_{B_n^i} \mu_n)_{n\in\sigma(\N)}$ converges weakly to $\mu^i$ by \labelcref{bubbleWeakConvergence}, and together with \labelcref{inductionBubbles1} it implies that\vspace*{-3pt}
\[(c_{B_n^i} \mu_n)_{n\in\sigma(\N)} \narrowto \mu^i,\]
\ie \labelcref{bubbleEmergence} is satisfied.
Moreover, by \labelcref{inductionBubbles1} again,\vspace*{-3pt}
\[
\limsup_{\ell\to\infty}\sum_{0\le i<k_{\sigma(\ell)}}\mu_{\sigma(\ell)}(B^i_{\sigma(\ell)})\le \sum_{0\le i<k}\norm{\mu^i}+\limsup_{\ell\to\infty}(\ell \wedge k) 2^{-\ell}=\sum_{0\le i<k}\norm{\mu^i},
\]
and since $k_n \to k$, by Fatou's lemma we have,\vspace*{-3pt}
\[\sum_{0\leq i < k}\norm{\mu^i} \leq \liminf_{\ell\to\infty}\sum_{0\le i<k_{\sigma(\ell)}}\mu_{\sigma(\ell)}(B^i_{\sigma(\ell)}),\]
which proves \labelcref{bubblingMassConservation} because $ \sum_{0\le i<k_{\sigma(\ell)}}\mu_{\sigma(\ell)}(B^i_{\sigma(\ell)}) = \norm{\mu^b_{\sigma(\ell)}}$.

\subsubsection*{Step 5: vanishing of the remaining part, proof of \labelcref{remainingVanishing}}
By \Cref{concentrationProperties}, it suffices to prove that \(m((\mu^v_n)_{n\in\sigma(\N)})=0\). We claim that:
\begin{equation}\label{bubbleConcentrationEstimate}
m((\mu^v_n)_{n\in\sigma(\N)})\le m((\mu^i_n)_{n\in\sigma_i(\N)}),\quad\text{for every \(i\in\N\) with \(i< k\)},
\end{equation}
which concludes since \(m((\mu^k_n)_{n\in\sigma_{k-1}(\N)})=0\) if \(k<\infty\), and \(m((\mu^i_n))_{n\in\sigma_i(\N)})\to 0\) as \(i\to\infty\) if \(k=\infty\). Indeed, if \(k=\infty\), we have by \labelcref{bubbleInduction4} and \labelcref{bubblingMassConservation},
\[
\frac 12\sum_{i\in\N}m((\mu^i_n)_{n\in\sigma_i(\N)})\le\sum_{i\in\N}\norm{\mu^i}= \lim_{\ell\to\infty} \norm{\mu_{\sigma(\ell)}^b} \leq \liminf_{\ell\to\infty}\norm{\mu_{\sigma(\ell)}}<\infty.
\]

Let us show \labelcref{bubbleConcentrationEstimate}. Let \(\bar\sigma\preceq\sigma\) and \((x_n)_{n\in\bar\sigma(\N)}\) be a sequence of points such that
\[
(\tau_{-x_n}\mu^v_n)_{n\in\bar\sigma(\N)}\mweakto \mu\in\Mspace(\R^N).
\]
We need to prove that \(\norm{\mu}\le m((\mu^i_n)_{n\in \sigma_i(\N)})\) for every $i <k$. Assume without loss of generality that \(\norm{\mu}>0\). Then for every \(i<k\),
\begin{equation}\label{distanceDiverge}
(\dist(x_n,x^i_n))_{n\in\bar\sigma(\N)}\to\infty.
\end{equation}
Otherwise, up to subsequence, $(\dist(x_n,x_n^i))_n$ would be bounded by some constant~$M$, and for every $r > 0$,
\[(\tau_{-x_n} \mu_n^v)(B_r) \leq \mu_n^v(B_{r+M}(x^i_n)) \xto{n\to\infty} 0,\]
because \(\mu^v_n\) is supported on \(\R^N\setminus\oldbigcup_{0\le i<k_n}B^i_n\) and $B_{r+M}(x^i_n) \subseteq B^i_n$ for $n$ large enough by \labelcref{bubbleDivergence}. Hence \(\mu\) would be \(0\), a contradiction. Up to further extraction, one can assume that \((\tau_{-x_n}\mu_n)_{n\in\bar\sigma(\N)}\) converges weakly to a measure \(\bar\mu \in \Mspace(\R^N)\). Since \(\mu^v_n\le\mu_n\), we~have \(\mu\le\bar\mu\). Moreover by \labelcref{bubbleInduction2}, for every \(i < k\) and \(n\in\ \bar\sigma(\N)\) large enough,
\[
\tau_{-x_n}\mu^i_n=\tau_{-x_n}\mu_n-\sum_{0\le j\leq i}\tau_{x^j_n-x_n}\mu^j,
\]
and because of \labelcref{distanceDiverge} the sum converges weakly to $0$, so that $\tau_{-x_n} \mu_n^i \mweakto \bar \mu$, and consequently,
\[
\norm{\mu}\le\norm{\bar\mu}\le m((\mu^i_n)_{n\in\sigma_i(\N)}),
\]
which is what had to be proved.

\subsubsection*{Step 6: re-centering of the bubbles at points of \(\supp\mu_n\)}
By \labelcref{bubbleWeakConvergence}, $(\tau_{-x^i_{n}} \mu_n)_{n \in \sigma(\N)}$ converges weakly to the non-trivial measure $\mu_i$ for every $i<k$, thus
\begin{equation}
\label{finiteDistance}
R_i/2\coloneqq \limsup_{\sigma(\N) \ni n \to +\infty} \dist(\supp \mu_n, x_n^i) < +\infty.
\end{equation}
Therefore, for every $n$ large enough, there is a point $\tilde x_n^i$ such that $\abs{x_n^i-\tilde x_n^i} < R_i$ and $\tilde x_n^i \in \supp \mu_n$. After a further extraction, one may assume that for every $i$, $\abs{x_n^i-\tilde x_n^i} < R_i < r_i^n$ where $\diam B_n^i = 2r_n^i$ for every $n$, and $(x_n^i -\tilde x_n^i)_{n\in \sigma(\N)}$ converges to some $p_i \in \R^N$. Finally, we set $\tilde r_i^n \coloneqq r_i^n-R_i$ and $\tilde B_n^i \coloneqq B(\tilde x_n^i,\tilde r_i^n) \subseteq B_n^i$. After replacing the balls $B_n^i$ by $\tilde B_n^i$, \labelcref{bubbleSplitting} and \labelcref{bubbleDivergence} are satisfied by definition. Notice that $(\tau_{-\tilde x^i_{n}} \mu_n)_{n \in \sigma(\N)}$ converges weakly to $\tilde \mu^i \coloneqq \tau_{p_i} \mu^i$ with $\norm{\tilde \mu^i} = \norm{\mu^i}$, and
\[
\limsup_n \norm{c_{B_n^i} \mu_n} = \limsup_n \mu_n(\tilde B_n^i) \leq \limsup_n \mu_n(B_n^i) = \norm{\mu^i}
\]
hence \labelcref{bubbleEmergence} holds. Besides, using Fatou's lemma,
\begin{align*}
\limsup_n \sum_{i<k_n} \mu_n(\tilde B_n^i) &\leq \limsup_n \sum_{i<k_n} \mu_n(B_n^i)\\
&= \sum_{i<k} \norm{\mu^i}\leq \sum_{i<k} \liminf_n \mu_n(\tilde B_n^i) \leq \liminf_n \sum_{i<k_n} \mu_n(\tilde B_n^i)
\end{align*}
so that $\lim_n \sum_{i<k_n} \mu_n(\tilde B_n^i) = \sum_i \norm{\mu_i}$ and \labelcref{bubblingMassConservation} is satisfied. In particular,
\[\lim_n \sum_{i<k_n} \mu_n(B_n^i\setminus \tilde B_n^i) = \lim_n \sum_{i<k_n} \mu_n(B_n^i) - \lim_n \sum_{i<k_n} \mu_n(\tilde B_n^i) = 0,\]
and \labelcref{remainingVanishing} holds as well.
\end{proof}
% Actualoriginal selected reference, officialpublishedphysicalPDF43/printed472
\begin{thebibliography}{Mar14}
\bibitem[Mar14]{marisProfileDecompositionSequences2014}
M.~Mari\c{s}, \emph{Profile decomposition for sequences of Borel measures}, 2014,
arXiv: \href{https://arxiv.org/abs/1410.6125}{1410.6125}.
\end{thebibliography}
\end{document}
